% SHARD-ID: RC-08B-WEIL-POSITIVITY
% SHARD-TITLE: Weil Positivity for an Arbitrary Transfer Operator
% SHARD-SUMMARY: Proves, for any finite-dimensional transfer operator, that positive definiteness of the rescaled trace sequence is exactly the one-sided bound on the retained spectrum, and that a duality of the spectrum turns it into "every mode is its own partner" and the Weil form into a mode-pairing form (the inflow identity).
% SHARD-SUMMARY: Sorts Kraus families by which half of RH each pairing buys: any Ad-family has positive rings and a conjugation symmetry; inverse pairing gives the functional equation; adjoint pairing gives reality of the transfer spectrum; both together is unitarity, where the criterion is Hastings' bound.
% SHARD-SUMMARY: Continuous (cMPS/Lindblad) version via Bochner with the Dyson-expansion ring norms; the Hilbert-Polya inner product exists iff semisimple on the circle; Weil positivity is blind to Jordan blocks; the zeta instance as a conditional dictionary; prior art (Weil, Li, Huang, Suzuki).
% SHARD-KEYWORDS: Weil positivity, explicit formula, inflow identity, functional equation, mode pairing, Herglotz, Bochner, Toeplitz, Hastings bound, inverse pairing, Hilbert-Polya, Jordan block, Li criterion, Huang criterion
% SHARD-NOTES: notes/weil-positivity.md
% SHARD-SCRIPTS: scripts/weil_positivity.py

\section{Weil Positivity for an Arbitrary Transfer Operator}
\label{sec:weil-positivity}

This section answers the question TJO put on 2026-09-12: the reflection
$\zero\mapsto1-\bar\zero$ pairs the mode of the Riemann channel at a zero with the
mode at its mirror image, RH says every mode is its own partner, and Weil's
criterion is the positivity of that pairing written on the primes. What does this mean
for a completely arbitrary Kraus family, finite-dimensional or continuous? The answer is
in three layers. Without any duality, Weil positivity of the rescaled ring norms is
exactly a one-sided bound, the spectral-gap half of RH. A duality of the spectrum, the
functional equation, turns the bound into a circle and the Weil form into the pairing of
each mode with its reflected partner; this identity is the finite form of the inflow
identity. For Kraus families the two pairings of the Kraus index buy the two halves
separately: inverse pairing gives the functional equation, adjoint pairing gives reality
of the transfer spectrum, and only unitarity gives both.

\emph{Status.} The theorems are the prover's, \texttt{codex:gpt-6-astra}
(\texttt{notes/weil-positivity/astra-proofs.md}, Lamport structure, correction ledger in
its T7), reviewed adversarially by \texttt{claude:opus}
(\texttt{notes/reviews/weil-positivity-2026-09-12.md}); two model families; all verdicts
\texttt{VALID}, two wording items fixed. The orchestrator's two observations in
\cref{sec:weil-positivity-continuous} were reviewed in the same file. Independent
numerics by a third agent (\cref{num:weil-positivity}) found the same three drafting
errors the prover corrected.

\subsection{Setting}

\Cref{def:rescaled-trace-sequence,def:weil-form} fix the objects: an operator $X$ on a
finite-dimensional space, a trivial sub-multiset $\trivS$ of its spectrum, the retained
multiset $\retA$, a critical radius $\critr$, the rescaled trace sequence
$\nuseq{\wl} = \sum_{\edgeev\in\retA}(\edgeev/\critr)^{\wl}$ with $\nuseq{-\wl} =
\overline{\nuseq{\wl}}$, the Weil form $\Wform(c) = \sum c_{\wl}\bar c_k\nuseq{\wl-k}$, the
reflection $\reflJ(\edgeev) = \critr^2/\bar\edgeev$ and the mode-pairing form $\Mform$. In
the Kraus case $X = \Hash$ and $\Tr\Hash^{\wl} = \sum_w|\Tr\Kr{w}|^2$ is the ring-norm sum
of \cref{cor:qihara-kraus}, so $\nuseq{\wl}$ is side A rescaled and with the trivial modes
subtracted, and $\Wform$ is a form on rings. Nothing below assumes diagonalisability
until \cref{prop:hp-inner-product-discrete}.

\subsection{Positivity is the one-sided bound}

\begin{theorem}[Finite Weil positivity]
\label{thm:weil-positivity-finite}
\claimstatus{thm:weil-positivity-finite}{proved}
For any $X$, any $\trivS$ and any $\critr>0$, the sequence $\nuseq{\wl}$ is positive definite
if and only if $|\edgeev|\le\critr$ for every $\edgeev\in\retA$, algebraic multiplicities
included and zero eigenvalues allowed.
\end{theorem}

The proof is the Poisson kernel and a growth lemma. A retained $z = \edgeev/\critr$
with $|z|<1$ contributes to $\Wform(c)$ the integral of $|\fcpoly{c}(e^{i\vartheta})|^2$
against the Poisson kernel $(1-|z|^2)/|1-ze^{-i\vartheta}|^2>0$, a point $|z| = 1$
contributes $|\fcpoly{c}(z)|^2$, and $z = 0$ contributes $\sum|c_{\wl}|^2$; so the bound
gives positivity. Conversely a positive definite sequence has $|\nuseq{\wl}|\le\nuseq{0}$
from its $2\times2$ minors, while a retained mode outside the disc makes the sequence
unbounded, because the modes of maximal modulus $R>\critr$ contribute $R^{\wl}$ times a
trigonometric sum whose mean square over $\wl<L$ tends to the sum of squared
multiplicities. The point of the theorem is what it does not say: an eigenvalue inside the
disc passes. Weil positivity alone is a spectral-gap statement, and the circle needs a
second input.

\subsection{A duality makes it two-sided: the inflow identity}

\begin{theorem}[Reflection duality and the inflow identity]
\label{thm:weil-duality-pairing}
\claimstatus{thm:weil-duality-pairing}{proved}
If $\retA$ is invariant under $\reflJ$ (multiplicities preserved; in particular
$0\notin\retA$), then $\Wform(c) = \Mform(c)$ for every finitely supported $c$, and the
following are equivalent: $\Wform\ge0$; $|\edgeev| = \critr$ for every $\edgeev\in\retA$;
$\reflJ(\edgeev) = \edgeev$ for every $\edgeev\in\retA$.
\end{theorem}

The identity $\Wform = \Mform$ is the finite inflow identity: the left side is built
from traces, i.e.\ from rings, the right side pairs each mode with its partner under the
reflection. It holds whether or not the modes are on the circle; what the circle adds is
that every partner is the mode itself, so $\Mform(c) = \sum|\fcpoly{c}(\edgeev/\critr)|^2$.
The equivalence follows from \cref{thm:weil-positivity-finite}: $|\edgeev|\le\critr$ and
$|\reflJ(\edgeev)|\le\critr$ force $|\edgeev| = \critr$. This is exactly how RH follows
from the one-sided bound $\re\zero\le\frac12$ together with the functional equation.

\begin{proposition}[Orbit contributions; an off-circle pair makes the form negative]
\label{prop:weil-orbit-negative}
\claimstatus{prop:weil-orbit-negative}{proved}
Under the hypotheses of \cref{thm:weil-duality-pairing}, a $\reflJ$-fixed mode $z$ of
multiplicity $m$ contributes $m|\fcpoly{c}(z)|^2$ to $\Wform(c)$ and a two-element orbit
$\{z,\reflJ z\}$ contributes $2m\,\re\big(\fcpoly{c}(z)\overline{\fcpoly{c}(\reflJ z)}\big)$;
if any two-element orbit exists there is a $c$ with $\Wform(c)<0$.
\end{proposition}

This is Weil's off-line pair argument in finite form; the prover's proof interpolates
$\fcpoly{c}$ to vanish at every other retained value, which is what makes the single
negative orbit dominate the whole form.

\begin{proposition}[Reciprocal symmetry without conjugation]
\label{prop:weil-reciprocal-only}
\claimstatus{prop:weil-reciprocal-only}{proved}
If $\retA$ contains no zero and is invariant under the reciprocal map
$\edgeev\mapsto\critr^2/\edgeev$ alone, positivity is still equivalent to $|\edgeev| =
\critr$ throughout $\retA$. If $\retA$ is also closed under complex conjugation it is
$\reflJ$-invariant and $\Wform = \Mform$; without conjugation symmetry $\Wform = \Mform$
can fail. The reciprocal map is holomorphic on $\mathbb C^\times$, not linear.
\end{proposition}

\subsection{Kraus families: which pairing buys which half}

\begin{proposition}[Conjugation symmetry needs no pairing]
\label{prop:kraus-conjugation-symmetry}
\claimstatus{prop:kraus-conjugation-symmetry}{proved}
For any family $\Eop{i} = \Ad(B_i)$ with arbitrary $B_i\in\Mn$ and any fixed-point-free
reversal, the antiunitary involution $x\mapsto x^\dagger$ on $\Vsp = \Mn$, and its
componentwise version on $\Wsp$, commute with $\Sig$ and with $\Hash$. Hence
$\operatorname{spec}(\Sig)$ and $\operatorname{spec}(\Hash)$ are closed under complex
conjugation with multiplicities. The nonnegative ring trace $\Tr\Hash^{\wl} =
\sum_w|\Tr B_w|^2$ of \cref{cor:qihara-kraus} likewise holds for every such family, with
no pairing hypothesis: its proof uses only $\Ad(A)\Ad(B) = \Ad(AB)$ and $\Tr\Ad(A) =
|\Tr A|^2$.
\end{proposition}

The draft sent to the prover asked for an inverse-paired counterexample to conjugation
closure; none exists inside the Kraus class, and the independent numerics found the same.
Likewise the trace formula $\Tr\Hash^{\wl} = \sum_w|\Tr B_w|^2$ of \cref{cor:qihara-kraus}
never used the pairing: it is a property of the $\Ad$ form. So every Kraus-type family,
however paired, has a positive side A and a real-structured side B.

\begin{theorem}[Inverse pairing gives the functional equation]
\label{thm:kraus-inverse-pairing-duality}
\claimstatus{thm:kraus-inverse-pairing-duality}{proved}
Let $\Eop{\bari{i}} = \Eop{i}^{-1}$ on an $N$-dimensional space, $\Dk = 2\mpairs\ge4$,
$\qs = \Dk-1$, $k = N(\Dk-2)/2$. Then $\operatorname{spec}(\Hash)$ is the multiset
$\{+1^{[k]},-1^{[k]}\}$ together with, for each eigenvalue $\sigev$ of $\Sig$ of
multiplicity $a_\sigev$, the two roots of $\edgeev^2-\sigev\edgeev+\qs$ counted $a_\sigev$
times. The retained multiset after removing $\{+1^{[k]},-1^{[k]}\}$ has $2N$ elements,
no zero, and is invariant under $\edgeev\mapsto\qs/\edgeev$; it is conjugation-invariant
iff $\operatorname{spec}(\Sig)$ is. For inverse-paired Kraus matrices both hold, so the
retained multiset is $\reflJ$-invariant for $\critr = \sqrt{\qs}$.
\end{theorem}

The bookkeeping is where the draft was loose: a double root at $\sigev = \pm2\sqrt{\qs}$
gives $2a_\sigev$ copies of $\pm\sqrt{\qs}$, and $\sigev = \pm\Dk$ gives extra copies of
$\pm1$ beyond the $k$ from the $(1-\uz^2)$ factor.

\begin{proposition}[Both pairings at once is unitarity]
\label{prop:kraus-both-pairings-unitary}
\claimstatus{prop:kraus-both-pairings-unitary}{proved}
$\Ad(B^\dagger) = \Ad(B)^{-1}$ iff $B^\dagger B = 1$ iff $B$ is unitary; a family is
simultaneously adjoint-paired and inverse-paired iff every $B_i$ is unitary. A scalar
$B^\dagger B = \kappa1$ gives only $\Ad(B^\dagger) = \kappa^2\Ad(B)^{-1}$; $B = 2$ disproves
``unitary up to a scalar''.
\end{proposition}

\begin{theorem}[The Weil criterion for Kraus families]
\label{thm:kraus-weil-criterion}
\claimstatus{thm:kraus-weil-criterion}{proved}
(i) For an adjoint-paired Kraus family, any $\trivS$ and any $\critr$, Weil positivity of
$\nuseq{\wl}(\Hash;\trivS,\critr)$ is exactly the one-sided bound $|\edgeev|\le\critr$ off
$\trivS$; nonnegativity of the ring traces alone does not give it. (ii) For unitary,
adjoint-paired $B_i$ with $\Dk\ge4$, $\critr = \sqrt{\Dk-1}$, $\chan = \Sig/\Dk$, and
$\trivS = \{+1^{[k]},-1^{[k]}\}$ together with $\{\qs,1\}$ once for each eigenvalue $+1$
of $\chan$ and $\{-\qs,-1\}$ once for each eigenvalue $-1$: Weil positivity holds iff every
other eigenvalue $\lambda$ of $\chan$ satisfies $|\lambda|\le2\sqrt{\Dk-1}/\Dk$ iff every
retained mode is $\reflJ$-fixed. Here $\lambda$ is real because $\chan$ is Hilbert--Schmidt
self-adjoint.
\end{theorem}

Part (ii) is Hastings' bound (\cref{cit:hastings-bound}) recovered as Weil positivity, and
it is the quantum-channel case of \cref{obs:rh-iff-ramanujan-channel}: the Ramanujan
property is the statement that the rescaled ring norms form a positive definite sequence.

\begin{proposition}[An adjoint-paired family with no duality]
\label{prop:kraus-no-duality-example}
\claimstatus{prop:kraus-no-duality-example}{proved}
For $n = 2$, $\Dk = 4$, $B_1 = B_3 = \operatorname{diag}(1,2)$, $B_2 = B_4 = 1$, the
retained spectrum of $\Hash$ is invariant under $\edgeev\mapsto c/\edgeev$ for no nonzero
$c$. The Weil form is nevertheless defined and its positivity is the one-sided bound.
\end{proposition}

So for a general channel ``every mode is its own partner'' has, literally, no partner to
refer to: the reflection exists on the punctured plane but need not send a mode to a mode.
What remains is \cref{thm:weil-positivity-finite}: a spectral-gap statement, the analogue
of $\re\zero\le\frac12$ without the functional equation.
